\chapter{内存与指针}
\label{chap:memory}

第~\ref{chap:c-basics}~章讲 \texttt{scanf} 时留了一个悬念：读入变量时，变量名前面为什么要加一个 \texttt{\&}？要回答这个问题，就得先弄清楚：程序里的变量，到底放在哪里？

这一章的前半部分讲内存：程序在内存里是什么样子，数据放在哪里，什么是栈，什么是堆，为什么用户程序碰不得操作系统内核。这些是计算机体系结构的基本概念，初学 C 的人不一定马上用得上；但心里有了这幅图，以后学指针、学操作系统、学 CPU，都会容易得多。后半部分讲指针——一个装着地址的变量。有了前面这幅图，指针就一点也不神秘了。

本章的示例在 \texttt{code/ch04/} 目录中。运行 \texttt{make} 编译，\texttt{make test} 检查。程序会打印出变量的地址，这些地址每次运行都不一样（原因见第~\ref{sec:layout}~节），书中的结果是在关闭了地址随机化的情况下得到的。你那里看到的数，多半和书上不同，这很正常。

\section{内存：一排有门牌号的房子}
\label{sec:mem-street}

内存可以想象成一条很长很长的街，街边是一排一模一样的小房子。每座房子里只能住一个\emph{字节}（byte），也就是 8 个二进制位，能表示 0 到 255 之间的一个数。每座房子都有一个门牌号，这就是\emph{地址}。CPU 要读写内存，就得说出门牌号：“把 0x7fffffffe560 号房子里的东西拿给我。”

在 C 语言里，变量名前面加上 \texttt{\&}，就能得到这个变量的地址：

\codefile{code/ch04/addr.c}{变量住在哪里}{lst:addr}

\inputminted[linenos=false]{text}{code/ch04/out/addr.txt}

几点说明：
\begin{itemize}
  \item \texttt{\&a} 读作“a 的地址”。\texttt{printf} 用 \texttt{\%p} 打印地址。前面的 \texttt{(void *)} 是告诉 \texttt{printf}“这是一个地址”，照抄即可。
  \item \texttt{sizeof(int)} 算出一个 \texttt{int} 占几个字节，要用 \texttt{\%zu} 打印。在我们的电脑上，一个 \texttt{int} 占 4 个字节，也就是要住 4 座房子。
  \item 地址是用\emph{十六进制}写的，前面的 \texttt{0x} 就是十六进制的标志。十六进制逢十六进一，用 \texttt{0}--\texttt{9} 和 \texttt{a}--\texttt{f} 这 16 个符号，\texttt{a} 是 10，\texttt{f} 是 15，\texttt{0x10} 是 16。地址很大，用十进制写太长，用十六进制正好：两个十六进制数字恰好是一个字节。
\end{itemize}

\texttt{a} 住在 \texttt{0x7fffffffe560}，\texttt{b} 住在 \texttt{0x7fffffffe564}，正好相差 4：\texttt{a} 占了 \texttt{...e560} 到 \texttt{...e563} 这 4 座房子，\texttt{b} 紧挨着住在后面（图~\ref{fig:houses}）。

\begin{figure}[htbp]
  \centering
  \begin{tikzpicture}[line width=0.7pt]
    \foreach \i/\v/\n in {0/03/60, 1/00/61, 2/00/62, 3/00/63, 4/04/64, 5/00/65, 6/00/66, 7/00/67} {
      \pgfmathsetmacro\x{\i*1.3}
      \filldraw[fill=yellow!15] (\x,0) rectangle +(1.05,0.9);
      \filldraw[fill=red!45] (\x-0.1,0.9) -- +(0.625,0.5) -- +(1.25,0) -- cycle;
      \filldraw[fill=brown!40] (\x+0.38,0) rectangle +(0.3,0.35);
      \node[font=\ttfamily\small] at (\x+0.525,0.63) {\v};
      \node[font=\ttfamily\scriptsize] at (\x+0.525,-0.25) {...e5\n};
    }
    % 烟囱和炊烟
    \filldraw[fill=gray!40] (0.75,1.15) rectangle (0.92,1.55);
    \foreach \r/\dx/\dy in {0.08/0.05/0.12, 0.11/0.2/0.32, 0.14/0.42/0.58}
      \draw[gray!70] (0.83+\dx,1.55+\dy) circle (\r);
    \draw[decorate, decoration={brace, mirror, amplitude=6pt}]
      (0,-0.5) -- (5.0,-0.5) node[midway, below=7pt] {\texttt{int a = 3}，占 4 个字节};
    \draw[decorate, decoration={brace, mirror, amplitude=6pt}]
      (5.2,-0.5) -- (10.15,-0.5) node[midway, below=7pt] {\texttt{int b = 4}，占 4 个字节};
    \node[font=\small, anchor=west] at (10.4,0.6) {← 房子里住的字节};
    \node[font=\small, anchor=west] at (10.4,-0.25) {← 门牌号（地址）};
  \end{tikzpicture}
  \caption{内存是一排有门牌号的房子，每座房子住一个字节}
  \label{fig:houses}
\end{figure}

仔细看图，3 这个数是怎样放进 4 个字节里的？它存成了 \texttt{03 00 00 00}：数值最低的那个字节，放在地址最小的房子里。这种存法叫\emph{小端}（little-endian）。x86 和 RISC-V 的处理器都是这样存数的。

\section{程序住在内存的哪里}
\label{sec:layout}

一个程序运行起来，它的代码、全局变量、局部变量都要在内存里找地方住。它们住在哪里？写个程序看看：

\codefile{code/ch04/layout.c}{程序的各个部分住在哪里}{lst:layout}

程序的前半部分打印出各种东西的地址：
\inputminted[linenos=false]{text}{code/ch04/out/layout.txt}

后半部分请操作系统给出这个程序的“内存地图”。在 Linux 上，每个正在运行的程序都有一个特殊的文件 \texttt{/proc/self/maps}，里面列出了这个程序用到的每一段内存。为了排版，这里只保留了其中几行，每行也只留下了地址范围、权限和名字三列：
\inputminted[linenos=false]{text}{code/ch04/out/layout_maps.txt}

把两边对照一下，就能画出程序在内存里的样子（图~\ref{fig:layout}）。它像一栋楼：地址从下往上越来越大，不同的东西住在不同的楼层。

\begin{figure}[htbp]
  \centering
  \begin{tikzpicture}[
      floor/.style={draw, minimum width=46mm, anchor=south west, inner sep=0pt},
      addr/.style={font=\ttfamily\scriptsize, anchor=east},
      what/.style={font=\small, anchor=west},
      line width=0.7pt,
    ]
    % 楼层（从下往上）
    \node[floor, minimum height=5mm, fill=gray!25]   at (0,0)    {不能访问（地址 0 附近）};
    \node[floor, minimum height=6mm, fill=blue!18]   at (0,0.5)  {代码};
    \node[floor, minimum height=6mm, fill=blue!8]    at (0,1.1)  {只读数据};
    \node[floor, minimum height=6mm, fill=green!18]  at (0,1.7)  {数据};
    \node[floor, minimum height=6mm, fill=green!8]   at (0,2.3)  {BSS};
    \node[floor, minimum height=8mm, fill=orange!25] at (0,2.9)  {堆 \quad $\uparrow$};
    \draw[dashed] (0,3.7) rectangle (4.6,5.3);
    \node[font=\small, gray] at (2.3,4.0) {（大片空地）};
    \node[floor, minimum height=5mm, fill=violet!12] at (0,4.4)  {共享库（libc 等）};
    \node[floor, minimum height=8mm, fill=red!18]    at (0,5.3)  {栈 \quad $\downarrow$};
    % 铁丝网
    \draw[gray!70!black, line width=1.2pt]
      (0,6.25) \foreach \k in {1,...,23} { -- ++(0.1,{mod(\k,2) ? 0.12 : -0.12}) };
    \foreach \k in {0.3,1.1,1.9,2.7,3.5,4.3}
      \draw[gray!70!black, line width=1pt] (\k,6.15) -- ++(0.12,0.2) (\k+0.12,6.15) -- ++(-0.12,0.2);
    \node[floor, minimum height=14mm, fill=gray!45] at (0,6.45) {};
    \node[font=\bfseries] at (1.6,7.45) {内核空间};
    % 告示牌
    \node[draw, fill=yellow!60, font=\scriptsize\bfseries, rotate=-6] at (1.6,6.95) {闲人免进};
    % 门卫
    \draw (3.85,7.45) circle (0.13);
    \fill[blue!60!black] (3.7,7.55) rectangle (4.0,7.62);
    \fill[blue!60!black] (3.75,7.62) rectangle (3.95,7.7);
    \draw (3.85,7.32) -- (3.85,6.95);
    \draw (3.85,6.95) -- (3.72,6.6) (3.85,6.95) -- (3.98,6.6);
    \draw (3.85,7.22) -- (3.6,7.42);
    \draw (3.85,7.22) -- (4.1,7.05);
    % 屋顶
    \filldraw[fill=red!40] (-0.2,7.85) -- (2.3,8.6) -- (4.8,7.85) -- cycle;
    % 左边：地址
    \node[addr] at (-0.15,0.05) {0};
    \node[addr] at (-0.15,0.8)  {0x555555555229};
    \node[addr] at (-0.15,1.4)  {0x555555556020};
    \node[addr] at (-0.15,2.0)  {0x555555558010};
    \node[addr] at (-0.15,2.6)  {0x555555558040};
    \node[addr] at (-0.15,3.3)  {0x555555559010};
    \node[addr] at (-0.15,5.7)  {0x7fffffffe550};
    \node[addr] at (-0.15,6.5)  {0xffff800000000000};
    \node[addr] at (-0.15,7.8)  {0xffffffffffffffff};
    % 右边：住户和权限
    \node[what] at (4.75,0.8) {\texttt{main} \quad \texttt{r-x}};
    \node[what] at (4.75,1.4) {\texttt{"hello"} \quad \texttt{r--}};
    \node[what] at (4.75,2.0) {\texttt{counter} \quad \texttt{rw-}};
    \node[what] at (4.75,2.6) {\texttt{table} \quad \texttt{rw-}};
    \node[what] at (4.75,3.3) {\texttt{malloc} 借来的内存，往上长};
    \node[what] at (4.75,5.7) {局部变量 \texttt{local}，往下长};
    \node[what] at (4.75,7.15) {操作系统自己用，用户程序碰不得};
    % 高低地址
    \draw[-{Stealth}, thick] (-3.3,0) -- (-3.3,7.8) node[above] {高地址};
    \node[below] at (-3.3,0) {低地址};
  \end{tikzpicture}
  \caption{一个程序的内存布局（x86-64 Linux，地址来自代码~\ref{lst:layout} 的运行结果；楼层的高度不按比例）}
  \label{fig:layout}
\end{figure}

从下往上，各层住的是（表~\ref{tab:segments}）：

\begin{table}[htbp]
  \centering
  \small
  \caption{程序的各个部分}
  \label{tab:segments}
  \begin{tabular}{L{2.2cm}L{4.6cm}L{2.6cm}L{1.6cm}}
    \toprule
    楼层 & 住的是什么 & 例子 & 权限 \\
    \midrule
    代码段（text）   & 编译出来的机器指令 & \texttt{main} 函数 & \texttt{r-x} \\
    只读数据段       & 不能修改的常量，比如字符串 & \texttt{"hello"} & \texttt{r--} \\
    数据段（data）   & 有初值的全局变量 & \texttt{counter = 100} & \texttt{rw-} \\
    BSS 段          & 没有初值的全局变量，程序开始时全是 0 & \texttt{table[1000]} & \texttt{rw-} \\
    堆（heap）       & 程序运行时向系统借来的内存，从下往上长 & \texttt{malloc} 的结果 & \texttt{rw-} \\
    栈（stack）      & 函数的局部变量，从上往下长 & \texttt{local} & \texttt{rw-} \\
    \bottomrule
  \end{tabular}
\end{table}

内存地图的第二列是\emph{权限}：\texttt{r} 表示可以读，\texttt{w} 表示可以写，\texttt{x} 表示可以当作指令来执行，\texttt{-} 表示不可以；最后的 \texttt{p} 表示这段内存是这个程序私有的。注意代码段是 \texttt{r-x}：可以读、可以执行，但不能写，所以程序不能在运行时偷偷改掉自己的指令；数据段是 \texttt{rw-}：可以读写，但不能当指令执行，所以混进数据里的恶意代码也跑不起来。这些检查不是靠操作系统一条一条去看的，而是 CPU 在每一次访问内存时，用硬件自动完成的。

楼的最底下，地址 0 附近，是一片“不能访问”的区域，不属于任何一段。后面会看到，这对发现程序的错误很有用。

\paragraph{地址为什么每次都不一样？}
如果你直接运行 \texttt{./layout}，再运行一次，会发现地址变了。这是操作系统故意的：每次运行，都把各段放到随机的位置上。这叫\emph{地址空间布局随机化}（ASLR），坏人想攻击一个程序，往往要先知道某样东西的地址；地址每次都变，攻击就难得多。为了让书中的结果每次都一样，我们用下面的命令关掉了随机化：
\begin{minted}[linenos=false]{text}
setarch $(uname -m) -R ./layout
\end{minted}

\paragraph{两个程序，同一个地址？}
如果同时运行两个程序，它们打印出的地址可能一模一样。难道两个程序住在同一个地方？不是的。每个程序都以为自己独占了一整栋楼，它看到的地址，其实是\emph{虚拟地址}。CPU 和操作系统一起，在背后把每个程序的虚拟地址，悄悄地对应到真正的内存（\emph{物理地址}）上，各住各的，互不干扰。这就是\emph{虚拟内存}。它是怎样做到的，等我们自己造 CPU 时再细讲。

\section{栈：倒着摞的盘子}
\label{sec:stack}

\subsection{先认识函数}

要讲栈，得先认识\emph{函数}。函数就是一段起了名字的代码，就像 \texttt{main} 一样。我们也可以自己写：

\codefile{code/ch04/func.c}{自己写一个函数}{lst:func}

\inputminted[linenos=false]{text}{code/ch04/out/func.txt}

\texttt{int add(int x, int y)} 的意思是：这个函数叫 \texttt{add}，接受两个整数 \texttt{x} 和 \texttt{y}（叫做\emph{参数}），算完以后交回一个整数。\texttt{return sum;} 就是把 \texttt{sum} 的值交回去。\texttt{main} 里写 \texttt{add(3, 4)}，就是\emph{调用}这个函数：3 放进 \texttt{x}，4 放进 \texttt{y}，执行 \texttt{add} 的代码，交回来的 7 放进 \texttt{c}。

如果一个函数不需要交回任何东西，就把返回类型写成 \texttt{void}；如果不需要参数，括号里也写 \texttt{void}。函数要写在调用它的地方前面，编译器才认识它。

\subsection{函数调用和栈}

函数里的变量，比如 \texttt{add} 里的 \texttt{sum}，叫\emph{局部变量}。它们住在哪里？下面的程序让 \texttt{main} 调用 \texttt{first}，\texttt{first} 调用 \texttt{second}，\texttt{second} 又调用 \texttt{third}，每个函数都打印出自己的局部变量的地址：

\codefile{code/ch04/stack.c}{一层一层地调用}{lst:stack}

\inputminted[linenos=false]{text}{code/ch04/out/stack.txt}

调用得越深，地址越小，每一层都比上一层小 \texttt{0x20}，也就是 32 个字节。每调用一个函数，程序就在\emph{栈}上为它划出一块地方，叫做这个函数的\emph{栈帧}，用来放它的局部变量；函数返回时，这块地方就还回去了。

栈就像一摞盘子：每调用一个函数，就在最上面加一个盘子；函数返回，就把最上面的盘子拿走。后放上去的先拿走，这正是函数调用的规律——\texttt{third} 最后被调用，却最先返回。只不过，内存里的这摞盘子是倒着摞的：它从高地址开始，往低地址长（图~\ref{fig:plates}）。

\begin{figure}[htbp]
  \centering
  \begin{tikzpicture}[line width=0.7pt]
    % 天花板
    \fill[brown!30] (-2.6,0.15) rectangle (2.6,0.45);
    \draw (-2.6,0.15) -- (2.6,0.15);
    \foreach \x in {-2.4,-2.0,...,2.4} \draw[brown!60] (\x,0.45) -- ++(0.25,-0.3);
    % 倒着摞的盘子
    \foreach \i/\n/\a in {0/main/e564, 1/first/e544, 2/second/e524, 3/third/e504} {
      \pgfmathsetmacro\y{-0.35-\i*0.62}
      \filldraw[fill=white]  (0,\y) ellipse (1.7 and 0.3);
      \filldraw[fill=cyan!8] (0,\y) ellipse (1.15 and 0.19);
      \node[font=\ttfamily\small] at (0,\y) {\n};
      \node[font=\ttfamily\scriptsize, anchor=west] at (2.0,\y) {...\a};
    }
    \node[font=\small, anchor=west] at (2.0,0.6) {地址};
    % 往下仰望的小人
    \draw (-0.2,-3.35) circle (0.17);
    \draw (-0.2,-3.52) -- (-0.2,-4.05);
    \draw (-0.2,-4.05) -- (-0.38,-4.50) (-0.2,-4.05) -- (-0.02,-4.50);
    \draw (-0.2,-3.70) -- (-0.5,-3.45);
    \draw (-0.2,-3.70) -- (0.1,-3.45);
    \node[red!70!black, font=\Large\bfseries] at (0.45,-3.05) {?!};
    \draw[gray, line width=1pt] (-2.6,-4.50) -- (2.6,-4.50);
    % 说明
    \draw[-{Stealth}, thick] (-2.3,-0.4) -- node[left, align=right, font=\small] {调用一层，\\往下多一个盘子} (-2.3,-2.2);
  \end{tikzpicture}
  \caption{栈是一摞从天花板往下摞的盘子}
  \label{fig:plates}
\end{figure}

这件事有两个重要的后果：
\begin{itemize}
  \item \textbf{局部变量是“临时住户”。}函数一返回，它的盘子就被拿走了，局部变量也就不存在了，那块地方马上会被下一次函数调用占用。所以，千万不要把局部变量的地址带出函数，以后再去用它。
  \item \textbf{栈是有限的。}在我们的 Linux 上，一个程序的栈最多 8\,MB（用 \texttt{ulimit -s} 命令可以查到，单位是 KB）。如果函数没完没了地调用下去，盘子越摞越多，栈就会用完，程序崩溃。这叫\emph{栈溢出}（stack overflow）。程序员最常去的问答网站 Stack Overflow，名字就是这么来的。
\end{itemize}

\section{堆：借来的内存，要记得还}
\label{sec:heap}

栈用起来很方便，进入函数自动分配，离开函数自动回收。但它有两个限制：大小有限；而且局部变量活不过它的函数。如果事先不知道要多少内存（比如要存全班同学的分数，可是班上有几个人，运行时才知道），或者一块内存要在函数返回之后继续用，就要用\emph{堆}。

堆是一大片空地，程序可以在运行时向系统“借”：
\begin{itemize}
  \item \texttt{malloc(n)} 借 \texttt{n} 个字节，交回这块内存的地址；如果借不到，交回 \texttt{NULL}（一个特殊的、表示“哪里也不是”的地址）。
  \item \texttt{free(地址)} 把借来的内存还回去。
\end{itemize}
它们定义在 \texttt{stdlib.h} 里。要真正用起这块借来的内存，需要用到指针，第~\ref{sec:heap-pointer}~节会看到完整的例子。

和栈不同，堆上的内存不会自动还：借了多少，就要还多少。忘了 \texttt{free}，这块内存就一直占着，谁也用不了，这叫\emph{内存泄漏}（memory leak）。偶尔漏一点问题不大，程序结束时，操作系统会把它用过的内存全部收回。但浏览器、服务器这些要连续运行好几天的程序，如果一直在漏，内存就会被一点一点地吃光，像一只漏水的桶（图~\ref{fig:leak}）。

\begin{figure}[htbp]
  \centering
  \begin{tikzpicture}[line width=0.8pt]
    % 水龙头
    \fill[gray!50] (-0.3,3.2) rectangle (1.1,3.45);
    \fill[gray!50] (0.85,3.2) rectangle (1.05,2.85);
    \node[font=\ttfamily\small, anchor=south] at (0.4,3.45) {malloc};
    \foreach \y in {2.45,2.05} \fill[cyan!60] (0.95,\y) circle (0.06);
    % 桶
    \filldraw[fill=cyan!25] (-0.45,0.1) -- (2.25,0.1) -- (2.45,1.3) -- (-0.65,1.3) -- cycle;
    \draw (-0.85,2.1) -- (2.65,2.1) -- (2.15,0) -- (-0.35,0) -- cycle;
    \draw (-0.85,2.1) .. controls (-0.6,3.0) and (2.4,3.0) .. (2.65,2.1);
    \node[font=\small] at (0.9,1.7) {内存};
    % 哭脸
    \fill (0.55,0.85) circle (0.05) (1.25,0.85) circle (0.05);
    \draw (0.6,0.45) arc[start angle=160, end angle=20, radius=0.32];
    \fill[cyan!60] (0.5,0.7) ellipse (0.03 and 0.07);
    % 漏洞和水滴
    \foreach \x/\y in {-0.1/-0.25, 0.9/-0.35, 1.8/-0.25} {
      \fill[cyan!60] (\x,\y) circle (0.06);
      \fill[cyan!60] (\x+0.05,\y-0.35) circle (0.05);
    }
    \fill[cyan!20] (0.9,-1.05) ellipse (1.9 and 0.15);
    \node[font=\small, anchor=west, align=left] at (3.0,0.6)
      {只 \texttt{malloc}，不 \texttt{free}：\\内存一点一点地漏光};
  \end{tikzpicture}
  \caption{内存泄漏}
  \label{fig:leak}
\end{figure}

还有两种更糟的错误：一块内存还回去以后又接着用（use after free），或者同一块内存还了两次（double free）。它们的后果难以预料，程序可能当场崩溃，也可能过了很久才莫名其妙地出错。第~\ref{chap:c-structure}~章的示例代码用 AddressSanitizer 检查这类错误，本章的 \texttt{make test} 也是这样做的。

\section{用户空间和内核空间}
\label{sec:kernel-space}

再看一眼图~\ref{fig:layout}：楼的最上面几层，是操作系统\emph{内核}的地盘，中间还隔着一道铁丝网。在 x86-64 的 Linux 上，64 位的地址被一分为二（表~\ref{tab:split}）。

\begin{table}[htbp]
  \centering
  \caption{x86-64 Linux 的地址划分（使用 4 级页表时）}
  \label{tab:split}
  \begin{tabular}{lll}
    \toprule
    地址范围 & 归谁 & 大小 \\
    \midrule
    \texttt{0x0000000000000000}--\texttt{0x00007fffffffffff} & 用户空间：普通程序 & 128\,TB \\
    \texttt{0x0000800000000000}--\texttt{0xffff7fffffffffff} & 不能使用 & \\
    \texttt{0xffff800000000000}--\texttt{0xffffffffffffffff} & 内核空间：操作系统 & 128\,TB \\
    \bottomrule
  \end{tabular}
\end{table}

我们写的程序，住在下半部分的\emph{用户空间}；上半部分的\emph{内核空间}，住着操作系统内核的代码和数据。内核管着整台计算机：内存怎么分，哪个程序先运行，文件怎么存，键盘和屏幕怎么用，都是它说了算。要是随便哪个程序都能改内核的数据，一个出错的程序（或者一个坏人写的程序）就能让整个系统瘫痪。所以，内核空间必须保护起来。

保护是靠 CPU 的硬件完成的。CPU 有两种工作状态：运行内核代码时处于\emph{内核态}，运行普通程序时处于\emph{用户态}（x86 把它们叫做 ring 0 和 ring 3，RISC-V 叫做 S 模式和 U 模式）。内核的每一页内存都标着“只有内核态才能访问”。CPU 每访问一次内存都会检查：如果它正处于用户态，却要访问内核的内存，就立刻拒绝，并通知内核来处理。内核的处理办法很简单：把这个程序结束掉。

不信就试试：

\codefile{code/ch04/kernel.c}{往内核空间里写数}{lst:kernel}

\inputminted[linenos=false]{text}{code/ch04/out/kernel.txt}

程序打印出第一行以后，就在 \texttt{*p = 42;} 这一行被结束了，“写进去了！”永远也打印不出来。（\texttt{*p = 42} 是什么意思，下一节就讲。）shell 告诉我们程序遇到了 \texttt{Segmentation fault}，中文叫\emph{段错误}——访问了不该访问的内存。\texttt{echo \$?} 显示程序的退出状态是 139，也就是 $128 + 11$，11 是段错误这个信号（SIGSEGV）的编号。

\begin{figure}[htbp]
  \centering
  \begin{tikzpicture}[line width=0.8pt]
    % 内核的大门
    \filldraw[fill=gray!45] (4.2,0) rectangle (6.4,3.0);
    \filldraw[fill=brown!50] (4.75,0) rectangle (5.85,2.0);
    \fill (5.7,1.0) circle (0.05);
    \node[font=\bfseries] at (5.3,2.5) {内核空间};
    \node[draw, fill=yellow!60, font=\scriptsize\bfseries] at (5.3,1.55) {闲人免进};
    % 门卫
    \draw (3.4,1.8) circle (0.2);
    \fill[blue!60!black] (3.15,1.95) rectangle (3.65,2.05);
    \fill[blue!60!black] (3.25,2.05) rectangle (3.55,2.2);
    \draw (3.4,1.6) -- (3.4,0.75);
    \draw (3.4,0.75) -- (3.2,0) (3.4,0.75) -- (3.6,0);
    \draw (3.4,1.4) -- (3.7,1.1);
    \draw (3.4,1.4) -- (2.95,1.75);
    \draw (2.85,1.85) circle (0.09);
    % 程序小人，拿着写着地址的纸条
    \draw (0.6,1.6) circle (0.2);
    \draw (0.6,1.4) -- (0.6,0.65);
    \draw (0.6,0.65) -- (0.4,0) (0.6,0.65) -- (0.8,0);
    \draw (0.6,1.25) -- (0.3,0.95);
    \draw (0.6,1.25) -- (1.1,1.35);
    \node[draw, fill=white, font=\ttfamily\tiny, rotate=8, inner sep=2pt] at (1.55,1.45) {0xffffffff81000000};
    \node[font=\scriptsize\ttfamily] at (0.6,-0.25) {./kernel};
    % 对话框
    \node[draw, rounded corners, fill=red!10, font=\small\bfseries, align=center] (stop) at (2.6,3.0) {Segmentation fault！};
    \draw (stop.south) ++(0.35,0) -- (3.25,2.1);
    \draw[gray, line width=1pt] (-0.3,0) -- (6.8,0);
  \end{tikzpicture}
  \caption{用户程序想闯进内核空间}
  \label{fig:guard}
\end{figure}

那么，用户程序要用屏幕、读文件，又不能碰内核的东西，该怎么办？答案是：去“窗口”办理。就像在银行，你不能自己走进金库去拿钱，只能在柜台的窗口递上一张单子，由工作人员替你去取。这个窗口就是\emph{系统调用}。还记得第~\ref{sec:runtime}~节吗？\texttt{puts} 最后正是通过 \texttt{write} 系统调用，请内核把字送到屏幕上的。CPU 有专门的指令来做这件事（x86-64 是 \texttt{syscall}，RISC-V 是 \texttt{ecall}）：执行它时，CPU 从用户态切换到内核态，跳到内核事先指定好的入口；内核办完事，再切换回用户态，回到程序中继续执行。

用户态和内核态、内存的权限检查、虚拟地址到物理地址的转换——这些都是 CPU 在硬件上实现的。等我们自己造 CPU 时，它们都是要亲手做出来的东西。

\section{指针：一张写着地址的纸条}
\label{sec:pointer}

现在可以讲指针了。

\emph{指针}也是一个变量，只不过它装的不是普通的数，而是一个地址——就像一张纸条，上面写着某座房子的门牌号。

“指针”的英文是 pointer。在英语里，pointer 还是一种猎犬的名字（中文叫指示犬）：它发现猎物时会停下来，抬起一只前爪，用鼻子直直地指着猎物所在的方向（图~\ref{fig:dog}）。C 语言里的指针做的也是这件事：它自己不是那个东西，但它告诉你那个东西在哪里。

\begin{figure}[htbp]
  \centering
  \begin{tikzpicture}[line width=0.8pt]
    % 尾巴：伸得笔直
    \draw[line width=2pt, line cap=round] (-1.05,0.25) -- (-1.95,0.55);
    % 后腿
    \draw[line width=2.4pt, line cap=round] (-0.7,-0.2) -- (-0.85,-1.05);
    \draw[line width=2.4pt, line cap=round] (-0.45,-0.25) -- (-0.4,-1.05);
    % 前腿：一条站着，一条抬起来
    \draw[line width=2.4pt, line cap=round] (0.55,-0.25) -- (0.6,-1.05);
    \draw[line width=2.4pt, line cap=round] (0.8,-0.2) -- (1.05,-0.5) -- (0.9,-0.72);
    % 身体
    \filldraw[fill=white] (0,0.1) ellipse (1.1 and 0.42);
    \fill[brown!70] (-0.35,0.25) ellipse (0.3 and 0.18);
    \fill[brown!70] (0.3,-0.05) ellipse (0.2 and 0.13);
    % 脖子和头
    \filldraw[fill=white] (0.85,0.2) -- (1.15,0.75) -- (1.45,0.65) -- (1.05,0.05) -- cycle;
    \filldraw[fill=white] (1.4,0.75) circle (0.3);
    \filldraw[fill=white] (1.55,0.6) -- (2.05,0.68) -- (2.05,0.85) -- (1.6,0.95) -- cycle;
    \fill (2.05,0.77) circle (0.06);
    \fill (1.5,0.85) circle (0.04);
    \filldraw[fill=brown!70] (1.25,0.95) .. controls (1.05,0.85) and (1.0,0.45) .. (1.2,0.35) .. controls (1.3,0.55) .. (1.38,0.85) -- cycle;
    % 项圈和名牌
    \draw[red!70!black, line width=2pt] (1.0,0.62) -- (1.3,0.45);
    \draw (1.15,0.55) -- (0.75,1.0);
    \node[draw, fill=yellow!40, font=\ttfamily\scriptsize, inner sep=2pt, anchor=east] at (0.9,1.1) {p = 0x7fffffffe55c};
    % 指向
    \draw[dashed, -{Stealth}, red!70!black] (2.15,0.77) -- (3.85,0.77);
    % 盒子 a
    \filldraw[fill=brown!30] (3.95,0.3) rectangle (5.15,1.25);
    \node[font=\huge\bfseries] at (4.55,0.78) {3};
    \node[draw, fill=yellow!40, font=\ttfamily, rotate=6, inner sep=2pt] at (4.55,0.0) {a};
    \node[font=\ttfamily\scriptsize] at (4.55,-0.45) {住在 0x7fffffffe55c};
    \draw[gray, line width=1pt] (-2.1,-1.05) -- (5.5,-1.05);
  \end{tikzpicture}
  \caption{指针（pointer）就像一只指示犬：它自己不是猎物，但它指着猎物在哪里}
  \label{fig:dog}
\end{figure}

\codefile{code/ch04/pointer.c}{第一个指针}{lst:pointer}

\inputminted[linenos=false]{text}{code/ch04/out/pointer.txt}

和指针有关的写法只有三种（表~\ref{tab:pointer}）。

\begin{table}[htbp]
  \centering
  \caption{指针的三种写法}
  \label{tab:pointer}
  \begin{tabular}{L{2.4cm}L{3.4cm}L{6.2cm}}
    \toprule
    写法 & 读作 & 意思 \\
    \midrule
    \texttt{int *p;} & p 是一个指向 int 的指针 & 准备一个盒子 \texttt{p}，它装的是一个 \texttt{int} 的地址 \\
    \texttt{p = \&a;} & 把 a 的地址给 p & 现在 \texttt{p} 里装着 \texttt{a} 的地址，我们说“\texttt{p} 指向 \texttt{a}” \\
    \texttt{*p} & p 指向的东西 & 顺着 \texttt{p} 里的地址找过去，那个格子里的东西 \\
    \bottomrule
  \end{tabular}
\end{table}

所以，\texttt{*p} 就是 \texttt{a}。读 \texttt{*p}，读到的是 \texttt{a} 的值 3；往 \texttt{*p} 里写 10，就是把 10 放进 \texttt{a} 的格子，\texttt{a} 就变成了 10。\textbf{指针最根本的用处，就是往一个指定的地址里存数据，或者从那里取数据。}

有一点容易弄混：星号 \texttt{*} 在两个地方出现，意思不一样。在声明 \texttt{int *p;} 里，它表示“\texttt{p} 是一个指针”；在 \texttt{*p = 10;} 这样的语句里，它表示“\texttt{p} 指向的那个格子”。

回头再看上一节的 \texttt{kernel.c}，现在就能看懂了。下面这一行让 \texttt{p} 指向内核空间里的一个地址：
\begin{minted}[linenos=false]{c}
int *p = (int *)0xffffffff81000000;
\end{minted}
\texttt{*p = 42;} 往那里写 42——于是被 CPU 当场拦了下来。地址前面的 \texttt{(int *)} 是说“把这个数当作一个 \texttt{int} 的地址”，照抄即可。

\paragraph{空指针。}
如果一个指针暂时不指向任何东西，就把它设为 \texttt{NULL}。\texttt{NULL} 其实就是地址 0。还记得图~\ref{fig:layout} 最底下那层“不能访问”的区域吗？正是因为地址 0 附近不能访问，一旦程序不小心通过一个 \texttt{NULL} 指针去读写，就会立刻得到一个段错误，而不是悄悄地改坏了某个地方的数据。错误发现得越早，越好找。

\section{解开 \texorpdfstring{\NoCaseChange{scanf}}{scanf} 的谜：指针做参数}
\label{sec:swap}

下面这个程序想交换两个变量的值：

\codefile{code/ch04/swap.c}{交换两个数}{lst:swap}

\inputminted[linenos=false]{text}{code/ch04/out/swap.txt}

\texttt{swap\_wrong} 为什么换不了？因为调用函数时，交给它的是值的\emph{副本}。\texttt{swap\_wrong(a, b)} 把 3 和 4 复制一份，放进 \texttt{swap\_wrong} 自己盘子里的 \texttt{x} 和 \texttt{y}。它把 \texttt{x} 和 \texttt{y} 换了过来，可是 \texttt{main} 盘子里的 \texttt{a} 和 \texttt{b} 根本没动（图~\ref{fig:swap} 左）。函数一返回，它的盘子被拿走，换好的副本也跟着没了。

\texttt{swap} 收到的则是 \texttt{a} 和 \texttt{b} 的地址。它通过 \texttt{*x} 和 \texttt{*y}，顺着地址直接去改 \texttt{main} 盘子里的格子，所以换成了（图~\ref{fig:swap} 右）。

\begin{figure}[htbp]
  \centering
  \begin{tikzpicture}[
      cell/.style={draw, fill=brown!25, minimum width=11mm, minimum height=7mm, font=\ttfamily},
      plate/.style={draw, rounded corners=8pt, fill=cyan!6, minimum width=42mm, minimum height=12mm},
      vname/.style={font=\ttfamily\small, left=1mm},
      line width=0.7pt,
    ]
    \foreach \ox/\who/\xa/\yb/\px/\py/\note in {
        0/swap\_wrong/4/3/3/4/{换的是副本，\\\texttt{a}、\texttt{b} 没变},
        5.6/swap/\&a/\&b/4/3/{顺着地址直接改，\\换成了}} {
      \begin{scope}[xshift=\ox cm]
        \node[plate] (m) at (0,0) {};
        \node[font=\small, above=1mm] at (m.north) {\texttt{main} 的盘子};
        \node[cell] (a) at (-0.45,0) {\px};  \node[vname] at (a.west) {a};
        \node[cell] (b) at ( 1.25,0) {\py};  \node[vname] at (b.west) {b};
        \node[plate] (s) at (0,-2.4) {};
        \node[font=\small, below=1mm] at (s.south) {\texttt{\who} 的盘子};
        \node[cell] (x) at (-0.45,-2.4) {\xa};  \node[vname] at (x.west) {x};
        \node[cell] (y) at ( 1.25,-2.4) {\yb};  \node[vname] at (y.west) {y};
        \node[font=\small, align=center] at (0,-4.45) {\note};
      \end{scope}
    }
    % 左边：复制了一份值
    \draw[-{Stealth}, gray, dashed] (-0.45,-0.45) -- node[right, font=\scriptsize] {复制} (-0.45,-1.95);
    \draw[-{Stealth}, gray, dashed] ( 1.25,-0.45) -- node[right, font=\scriptsize] {复制} ( 1.25,-1.95);
    % 右边：x、y 里装的是地址，指向 main 盘子里的 a、b
    \fill[yellow!30] (4.6,-2.75) rectangle (5.7,-2.05)  (6.3,-2.75) rectangle (7.4,-2.05);
    \node[cell, fill=yellow!30] at (5.15,-2.4) {\&a};
    \node[cell, fill=yellow!30] at (6.85,-2.4) {\&b};
    \draw[-{Stealth}, red!70!black, thick] (5.15,-2.05) -- (5.15,-0.45);
    \draw[-{Stealth}, red!70!black, thick] (6.85,-2.05) -- (6.85,-0.45);
  \end{tikzpicture}
  \caption{传值和传地址}
  \label{fig:swap}
\end{figure}

现在，第~\ref{sec:scanf}~节留下的谜可以解开了。\texttt{scanf} 要把读到的数放进你的变量里。如果写成 \texttt{scanf("\%d", n)}，交给 \texttt{scanf} 的只是 \texttt{n} 的值的一个副本，它没法改动真正的 \texttt{n}；写成 \texttt{scanf("\%d", \&n)}，交给它的是 \texttt{n} 的地址，它才能顺着地址把数放进去。

一个函数想改动调用者的变量，就要拿到这个变量的地址。这是指针最常见的用法。

\section{指针加一：在内存里走路}
\label{sec:heap-pointer}

现在可以把堆真正用起来了。下面的程序先问有几个同学，再向系统借一块刚好够用的内存，存下每个同学的分数：

\codefile{code/ch04/heap.c}{用堆上的内存存分数}{lst:heap}

\inputminted[linenos=false]{text}{code/ch04/out/heap.txt}

\begin{itemize}
  \item \texttt{malloc(n * sizeof(int))} 借了 \texttt{n} 个 \texttt{int} 那么大的一块内存，把它的地址放进指针 \texttt{scores}。借完要检查一下是不是 \texttt{NULL}。
  \item \texttt{scores + i} 是第 \texttt{i} 个格子的地址，\texttt{*(scores + i)} 就是那个格子里的分数。读入时直接把地址 \texttt{scores + i} 交给 \texttt{scanf}，不用再加 \texttt{\&}，因为它本来就是地址了。
  \item 用完以后，\texttt{free(scores)} 把内存还回去。注意中途出错提前返回时，也要先 \texttt{free}。
\end{itemize}

仔细看打印出来的地址：每往后一个，地址加 4，而不是加 1。这是因为 \texttt{scores} 是指向 \texttt{int} 的指针，C 语言知道一个 \texttt{int} 占 4 个字节，所以“指针加一”的意思是“走到下一个 \texttt{int}”，地址实际加的是 4。如果是指向 \texttt{char}（占 1 个字节）的指针，加一就真的是地址加 1。

\texttt{*(scores + i)} 也可以写成 \texttt{scores[i]}，两种写法完全一样。后一种写法叫\emph{数组}，以后会经常用到。

\section{往指定的地址写数据}
\label{sec:mmio}

指针说到底，就是“往一个指定的地址写数据，或者从那里读数据”。在 Linux 上，用户程序只能读写操作系统分给它的那些地址，写错了地方就是段错误。可是，如果没有操作系统呢？

我们要造的 CPU，一开始上面不会有操作系统。CPU 能访问的地址里，除了内存，还有各种\emph{设备}：串口、定时器、指示灯……每个设备都占着几个地址。往某个地址写一个数，可能就是往内存里存了一个数；也可能是往串口发送了一个字符，或者点亮了一盏灯。这种做法叫\emph{内存映射输入输出}（memory-mapped I/O）。比如，在 QEMU 模拟的 RISC-V virt 开发板上，串口的数据寄存器就在地址 \texttt{0x10000000}，往那里写一个字节，屏幕上就出现一个字符：
\begin{minted}[linenos=false]{c}
volatile unsigned char *uart = (volatile unsigned char *)0x10000000;

*uart = 'H';
*uart = 'i';
\end{minted}

\texttt{unsigned char} 是一个字节大小的类型，\texttt{'H'} 是字符 H 的编码。\texttt{volatile} 告诉编译器：每一次读写都要老老实实地真的去做，不许优化掉（第~\ref{chap:hello}~章说过，编译器很喜欢自作主张）。

这段代码不能在 Linux 上直接运行——那样只会得到一个段错误。但在没有操作系统的机器上，这就是“打印”的全部秘密。\texttt{printf} 绕了一大圈，经过 C 标准库、系统调用、操作系统的驱动程序，最后做的也正是这件事。等我们的 CPU 造好了，它说的第一句话，就会是这样写出来的。

\section{小结}

\begin{itemize}
  \item 内存是一排有门牌号的格子，每个格子一个字节，门牌号就是地址。\texttt{\&a} 是 \texttt{a} 的地址，用 \texttt{\%p} 打印。x86 和 RISC-V 都是小端存储。
  \item 程序在内存里分成代码、只读数据、数据、BSS、堆、栈几部分，各有各的读、写、执行权限，由 CPU 在硬件上检查。
  \item 栈存放函数的局部变量，每调用一个函数加一个栈帧，函数返回就收回。栈从高地址往低地址长，大小有限。
  \item 堆上的内存用 \texttt{malloc} 借、用 \texttt{free} 还，忘了还就是内存泄漏。
  \item 地址空间分为用户空间和内核空间。CPU 在用户态访问内核空间会被拒绝，程序得到段错误。用户程序通过系统调用请内核办事。
  \item 指针是装着地址的变量：\texttt{p = \&a} 让 \texttt{p} 指向 \texttt{a}，\texttt{*p} 就是 \texttt{a}。函数要改调用者的变量，就要拿到它的地址，这就是 \texttt{scanf} 要写 \texttt{\&} 的原因。
  \item 指针加一，走到下一个元素，地址加的是元素的大小。
  \item 在没有操作系统的机器上，往设备的地址写数据，就能控制设备。
\end{itemize}

\section*{练习}

\begin{enumerate}
  \item 在代码~\ref{lst:addr} 中再加一个 \texttt{char} 型的变量，打印它的地址和 \texttt{sizeof(char)}。
  \item 不加 \texttt{setarch}，把 \texttt{layout} 连续运行三次，看看哪些地址变了。再加上 \texttt{setarch \$(uname -m) -R} 运行三次，比较一下。
  \item 写一个函数 \texttt{void add\_one(int *p)}，把 \texttt{p} 指向的那个数加 1。在 \texttt{main} 中调用它，验证变量真的变了。
  \item 修改代码~\ref{lst:heap}，再打印出最高分。
  \item 想一想：如果一个函数把自己的局部变量的地址交回给调用者，调用者再通过这个地址去读，会出什么问题？
  \item 写一个程序，执行 \texttt{int *p = NULL; *p = 1;}，看看会发生什么。结合图~\ref{fig:layout}，解释原因。
\end{enumerate}
